% Laboratorio 02 — Ubuntu 24.04 LTS: estación receptora ADS-B/Mode S en 1090 MHz.

\PortadillaLaboratorio{Laboratorio 02 · Ubuntu 24.04 LTS · 225 minutos}{Del pulso de 1090 MHz a la aeronave en el mapa: recepción y decodificación ADS-B/Mode S}{Usted montará una estación receptora pasiva: el HackRF captura ráfagas de 1090 MHz de microsegundos, SDRangel busca preámbulos, comprueba la paridad de cada trama y reconstruye identificación, altitud, velocidad y posición. El objetivo no es «ver aviones», sino poder demostrar qué parte de lo que muestra la pantalla fue recibida, cuál fue calculada y cuál vino de una base de datos.}

\section{Ficha y planificación}
\label{sec:u-plan}

\subsection{Resultados de aprendizaje}

Al terminar, usted podrá:
\begin{enumerate}
  \item preparar Ubuntu 24.04 para recibir con un SDR;
  \item verificar el HackRF y sus permisos USB (reglas \emph{udev});
  \item explicar para qué sirven las transmisiones aeronáuticas en 1090 MHz;
  \item identificar las etapas de la decodificación ADS-B;
  \item configurar en SDRangel una cadena de recepción compatible;
  \item observar actividad en 1090 MHz cuando existan señales adecuadas;
  \item obtener tramas válidas en condiciones de recepción calificadas;
  \item interpretar una tabla de aeronaves;
  \item distinguir los campos de los distintos mensajes Mode S y ADS-B;
  \item mostrar en un mapa las aeronaves con posición válida;
  \item medir y discutir el rendimiento de la recepción;
  \item separar evidencia RF, información derivada y enriquecimiento externo.
\end{enumerate}

\subsection{Presupuesto de tiempo}

\begin{longtable}{@{}>{\raggedright\arraybackslash}p{0.12\linewidth}>{\raggedright\arraybackslash}p{0.51\linewidth}>{\raggedleft\arraybackslash}p{0.10\linewidth}>{\raggedleft\arraybackslash}p{0.14\linewidth}@{}}
\toprule
\textbf{Sección} & \textbf{Actividad} & \textbf{Minutos} & \textbf{Acumulado} \\
\midrule
\endhead
\ref{sec:u-diag}, \ref{sec:u-conceptos} & Diagnóstico; bloque A: sistema y arquitectura & 20 & 20 \\
\ref{sec:u-antena}, \ref{sec:u-ubuntu} & Bloques B y C: antena y comprobación de Ubuntu & 30 & 50 \\
\ref{sec:u-receptor}, \ref{sec:u-demod} & Bloques D y E: HackRF y demodulador ADS-B & 35 & 85 \\
\ref{sec:u-tramas} & Bloque F: de la energía a la trama válida (con recepción en curso) & 40 & 125 \\
\ref{sec:u-tabla}, \ref{sec:u-mapa} & Bloques G y H: tabla de aeronaves y mapa & 45 & 170 \\
\ref{sec:u-exp}, \ref{sec:u-siglent} & Bloques I y J: estadísticas y SIGLENT & 25 & 195 \\
\ref{sec:u-offline} & Bloque K: respaldo offline verificado (si hace falta) & 15 & 210 \\
\ref{sec:u-informe}, \ref{sec:u-final} & Bloque L: informe y cuestionario final & 15 & 225 \\
\bottomrule
\end{longtable}

\begin{importante}[Plan de respaldo, decidido al minuto 85]
Si al terminar el bloque E (minuto 85) no hay \emph{ninguna} trama válida tras 10 minutos de recepción con la antena en la mejor ubicación disponible, el grupo pasa al bloque K \emph{en ese momento} y hace los bloques F--H con la grabación. Los 15 minutos del bloque K ya están en el presupuesto; si hubo recepción en vivo, se usan para comparar ambas fuentes.
\end{importante}

\subsection{Material del puesto}

\begin{checklist}
  \item HackRF One, cable USB de datos, antena B y, para llegar a una ventana, cable C.
  \item PC con Ubuntu 24.04 LTS \emph{nativo} (una máquina virtual complica el USB sostenido), SDRangel 7.27.2 (snap) y paquete \texttt{hackrf}, instalados por el instructor.
  \item La grabación \texttt{datos\_adsb/modes1\_dump1090\_2MSps\_1090MHz.sdriq} copiada en el puesto.
  \item SIGLENT SSA3032X Plus con adaptador N (bloque J, condicional).
\end{checklist}

\section{Cuestionario diagnóstico}
\label{sec:u-diag}

Ocho minutos, individual, sin consultar.

\CuestionarioDiagnosticoU

\section{Bloque A --- El sistema: radar, Mode S y ADS-B}
\label{sec:u-conceptos}

\begin{longtable}{@{}>{\raggedright\arraybackslash}p{0.17\linewidth}>{\raggedright\arraybackslash}p{0.47\linewidth}>{\raggedright\arraybackslash}p{0.28\linewidth}@{}}
\toprule
\textbf{Sistema} & \textbf{Qué es} & \textbf{¿Lo observa esta estación?} \\
\midrule
\endhead
PSR & Radar primario: emite pulsos y recibe el eco reflejado. No necesita cooperación del avión; no identifica. & No. \\
SSR & Radar secundario: interroga en 1030 MHz y el transpondedor responde en 1090 MHz. & Solo las \emph{respuestas} en 1090 MHz. \\
Modo A/C & Respuestas clásicas: código de identidad de 4 dígitos (A) y altitud de presión (C). Sin dirección única. & No las decodifica este laboratorio. \\
Mode S & Interrogación selectiva por dirección ICAO de 24 bits; respuestas con paridad (DF0, 4, 5, 11, 16, 20, 21\ldots). & Sí: recibe respuestas a interrogaciones de \emph{otros}. \\
ADS-B Out & La aeronave difunde, sin ser interrogada, su posición (GNSS), velocidad e identificación. & Sí. \\
1090ES & \emph{Extended squitter}: el enlace ADS-B Out más usado; tramas largas DF17/DF18 de 112 bits. & Sí: es el núcleo del laboratorio. \\
ADS-B In / recepción & Recibir y decodificar esos mensajes, a bordo o en tierra. & Es lo que usted construye. \\
MLAT & Multilateración: posición a partir de diferencias de tiempo de llegada en \emph{varios} receptores sincronizados. & No: hace falta una red, no un receptor. \\
\bottomrule
\end{longtable}

\begin{alerta}[Por qué este laboratorio es pasivo]
Su estación \textbf{no interroga}. Las tramas Mode S que verá son respuestas a interrogaciones de radares secundarios cercanos o de equipos anticolisión (ACAS) de otras aeronaves, más las difusiones espontáneas ADS-B. Recibir ADS-B \textbf{no} equivale a operar un radar secundario: no elige a quién preguntar, no mide distancias por tiempo de ida y vuelta y no puede forzar una respuesta. Emitir en 1030 o 1090 MHz está prohibido (sección~\ref{sec:seguridad}).
\end{alerta}

\begin{figure}[H]
\centering
\resizebox{\linewidth}{!}{% Laboratorio 02: arquitectura física y lógica de la estación receptora pasiva.
\begin{tikzpicture}[node distance=5mm and 7mm]
  \node[bloqueext,minimum width=27mm] (xpdr) {Transpondedor de la\\aeronave (Mode S,\\ADS-B Out 1090ES)};
  \node[bloque,right=22mm of xpdr,minimum width=18mm] (ant) {Antena\\receptora};
  \node[bloque,right=of ant,minimum width=21mm] (hrf) {\textbf{HackRF One}\\1090,000 MHz\\(solo RX)};
  \node[bloque,right=of hrf,minimum width=20mm] (pc) {Ubuntu 24.04\\USB + udev};
  \node[bloquealt,right=of pc,minimum width=23mm] (in) {SDRangel\\\textbf{HackRF Input}\\I/Q a 4 MSa/s};
  \draw[flecha,decorate,decoration={snake,amplitude=1.1mm,segment length=3mm,post length=2.5mm}] (xpdr) -- (ant);
  \node[rotulo] at ($(xpdr.east)!0.5!(ant.west)+(0,6mm)$) {1090 MHz\\ráfagas PPM};
  \draw[flecha] (ant) -- (hrf);
  \draw[flecha] (hrf) -- (pc);
  \draw[flecha] (pc) -- (in);

  \node[bloquealt,below=11mm of in,minimum width=27mm] (dem) {\textbf{ADS-B Demodulator}\\preámbulo, bits, paridad};
  \node[bloquealt,left=of dem,minimum width=23mm] (msg) {Tramas válidas\\(DF, ICAO, TC)};
  \node[bloquealt,left=of msg,minimum width=22mm] (tab) {Tabla de\\aeronaves};
  \node[bloquealt,left=of tab,minimum width=24mm] (map) {Feature \textbf{Map}\\(posición CPR)};
  \draw[flecha] (in) -- (dem); \draw[flecha] (dem) -- (msg); \draw[flecha] (msg) -- (tab); \draw[flecha] (tab) -- (map);

  \node[bloqueext,below=10mm of map,minimum width=30mm] (pos) {\menu{Preferences} $\rightarrow$ \textbf{My Position}\\(referencia CPR y distancia)};
  \node[bloqueext,below=10mm of msg,minimum width=30mm,xshift=-6mm] (db) {Bases externas descargadas\\(enriquecen; \emph{no} son RF)};
  \node[bloqueinst,below=10mm of dem,minimum width=25mm,xshift=8mm] (file) {\textbf{FileInput} (.sdriq)\\respaldo grabado};
  \draw[flechadash] (pos) -- (map);
  \draw[flechadash] (db.north) -- (tab.south east);
  \draw[flechadash] (file.north) -- ([xshift=8mm]dem.south);
\end{tikzpicture}
}
\caption{Arquitectura de la estación. Las flechas continuas son el camino de la señal; las discontinuas, entradas que \emph{no} vienen de la aeronave: la posición del receptor, las bases de datos y, en el respaldo, el archivo grabado.}
\label{fig:cadenaadsb}
\end{figure}

\section{Bloque B --- Compatibilidad de antena a 1090 MHz}
\label{sec:u-antena}

\begin{prediccion}
Con la antena B, de banda probable 868/915 MHz y no calificada a 1090 MHz, ¿espera recibir aeronaves a 10 km? ¿A 100 km? ¿Qué factor cree que pesa más: la antena, su ubicación o el tráfico?
\end{prediccion}

Cuatro factores deciden la recepción a 1090 MHz:
\begin{itemize}
  \item \textbf{Banda de la antena:} fuera de su banda, la desadaptación devuelve parte de la potencia y el diagrama cambia. No hay especificación de la antena B a 1090 MHz.
  \item \textbf{Ubicación:} 1090 MHz se propaga casi en línea de vista. Paredes, vidrios metalizados y estructuras atenúan mucho; una ventana despejada vale más que varios dB de ganancia.
  \item \textbf{Ganancia y polarización:} las antenas de aeronave radian en polarización vertical; mantenga la antena B vertical.
  \item \textbf{Pérdidas del cable:} unos 2 dB en 3 m de coaxial delgado (sección~\ref{sec:inventario}); compensan si permiten llegar a la ventana.
\end{itemize}

\Paso{Registre la decisión de preparación de su antena}
{Tabla~\ref{tab:decision} y hoja de registro.}
{Conector (SMA/RP-SMA), ubicación (interior, ventana, exterior), con o sin cable C, orientación vertical.}
{Una configuración \COND\ solo pasa a calificada con evidencia de nivel 4.}
{Una línea en el informe: «Antena B, SMA, ventana al norte, cable C: \COND».}
{Tras el bloque E, actualice: «\PASS\ para este lugar: $N$ tramas válidas en 5 min» o «no superada».}
{RP-SMA sin adaptador: \BLOQ. Pase directamente al bloque K.}

\begin{nota}[Qué antena haría predecible la recepción en vivo]
Una antena especificada para 1085--1095 MHz (monopolo $\lambda/4$ con plano de tierra, o colineal), 50~$\Omega$, ROE $\le 2{:}1$ en esa banda, polarización vertical, conector SMA macho, instalada en exterior o en una ventana con cielo despejado y con un cable corto de bajas pérdidas. Es una recomendación de compra condicional, no equipo disponible.
\end{nota}

\section{Bloque C --- Preparación de Ubuntu 24.04}
\label{sec:u-ubuntu}

\begin{estado}
Órdenes y nombres de paquetes comprobados el 15-10-2026 en las fuentes oficiales: paquete \texttt{hackrf} 2023.01.1-9build1 (universo de Ubuntu 24.04, \emph{noble}), que instala la regla \texttt{/usr/lib/udev/rules.d/60-libhackrf0.rules}; SDRangel 7.27.2 en la Snap Store (editor f4exb). \textbf{No se ejecutaron en Ubuntu} durante esta revisión: el instructor las ensaya antes de clase.
\end{estado}

\begin{importante}[Por qué snap y no el .deb]
La release 7.27.2 publica \texttt{sdrangel\_7.27.2\_ubuntu-24.04\_amd64.deb}, pero ese paquete declara una dependencia de \texttt{libuhd4.10.0}, que no existe en los repositorios estándar de Ubuntu 24.04 (que ofrecen \texttt{libuhd4.6.0t64}). La instalación del \texttt{.deb} fallaría por dependencia no satisfecha sin repositorios adicionales. Por eso este manual usa el paquete snap oficial. Use una sola vía de instalación.
\end{importante}

Abra una terminal con \tecla{Ctrl}+\tecla{Alt}+\tecla{T}. Las órdenes con \texttt{sudo} piden su contraseña; lo que escribe no se muestra en pantalla.

\Comando{lsb_release -ds}
{Confirmar la versión del sistema.}
{\texttt{Ubuntu 24.04.x LTS}.}
{Otra versión (p.\ ej.\ 22.04): los nombres de paquetes pueden diferir.}
{Avise al instructor; no actualice el sistema durante la clase.}

\Comando{sudo apt update}
{Actualizar el índice de paquetes (no instala nada).}
{Líneas \texttt{Hit:}/\texttt{Get:} y, al final, \texttt{Reading package lists... Done}.}
{\texttt{Temporary failure resolving} (sin red) o un bloqueo de \texttt{dpkg}.}
{Sin red: el instructor ya instaló el software; continúe. Bloqueo: espere a que termine la actualización automática.}

\Comando{sudo apt install hackrf}
{Instalar las herramientas HackRF y la biblioteca \texttt{libhackrf0} con su regla udev. Normalmente ya está instalado (preparación previa).}
{\texttt{hackrf is already the newest version (2023.01.1-9build1)} o la instalación de ese paquete.}
{\texttt{Unable to locate package hackrf}: el repositorio \emph{universe} no está habilitado.}
{\texttt{sudo add-apt-repository universe}, luego repita \texttt{sudo apt update}. Tarea del instructor si no hay permisos.}

\Comando{lsusb -d 1d50:6089}
{Comprobar que el núcleo ve el HackRF en el bus USB (fabricante 1d50, producto 6089).}
{Una línea con \texttt{ID 1d50:6089} y una descripción del HackRF One.}
{Sin salida: el HackRF no está enumerado.}
{Otro cable de datos, otro puerto, sin concentrador. \texttt{sudo dmesg | tail} muestra los mensajes del núcleo al reconectar.}

\Comando{hackrf_info}
{Identificar el HackRF \emph{sin} \texttt{sudo}: así se comprueba que los permisos udev funcionan.}
{\texttt{Found HackRF}, \texttt{Board ID Number: 2 (HackRF One)}, número de serie y versión de firmware.}
{\texttt{No HackRF boards found} con \texttt{lsusb} correcto, o error de acceso: falta permiso. Otro proceso que usa el HackRF da también error.}
{Cierre SDRangel. Si funciona solo con \texttt{sudo hackrf\_info}, siga las dos órdenes siguientes. No ejecute SDRangel con \texttt{sudo}.}

\Comando{cat /usr/lib/udev/rules.d/60-libhackrf0.rules}
{Leer la regla que da acceso al HackRF a usuarios sin privilegios.}
{Líneas que contienen \texttt{1d50} y \texttt{6089} y fijan permisos (grupo, modo o etiqueta \texttt{uaccess}).}
{\texttt{No such file or directory}: el paquete \texttt{libhackrf0} no se instaló.}
{Reinstale con \texttt{sudo apt install --reinstall libhackrf0}.}

\Comando{sudo usermod -aG plugdev $USER}
{Solo si la regla asigna el grupo \texttt{plugdev} y la orden \texttt{groups} no lo muestra: añadir su usuario a ese grupo.}
{Sin salida (éxito silencioso).}
{\texttt{user ... does not exist}: escribió mal la orden.}
{Cierre la sesión de Ubuntu y vuelva a entrar (los grupos se cargan al iniciar sesión); desconecte y reconecte el HackRF; repita \texttt{hackrf\_info}.}

\Comando{sudo snap install sdrangel}
{Instalar SDRangel desde la Snap Store (preparación previa del instructor).}
{\texttt{sdrangel 7.27.2 from Edouard Griffiths (f4exb) installed} o \texttt{already installed}.}
{Error de red o de \texttt{snapd}.}
{Sin red en clase: el software debe venir instalado; avise.}

\Comando{sudo snap connect sdrangel:raw-usb}
{Conceder al snap acceso a dispositivos USB en bruto (necesario para el HackRF).}
{Sin salida (éxito silencioso).}
{Un mensaje que indique que el snap no tiene el \emph{plug} \texttt{raw-usb}: versión distinta de la verificada.}
{Compruebe con \texttt{snap connections sdrangel}: la fila \texttt{raw-usb} debe mostrar una conexión (no un guion).}

\Comando{snap run sdrangel}
{Arrancar SDRangel como usuario normal (también desde el menú de aplicaciones).}
{Se abre la ventana principal.}
{Error gráfico, o la ventana no aparece; el snap exige una CPU con SSE 4.2.}
{Ejecute desde la terminal para leer el mensaje de error y anótelo para el instructor.}

\begin{alerta}[Un solo programa por HackRF]
El HackRF solo admite una aplicación a la vez. Cierre SDRangel antes de \texttt{hackrf\_info} y viceversa. Nunca ejecute SDRangel con \texttt{sudo}: archivos de configuración creados como \emph{root} causan fallos posteriores difíciles de diagnosticar.
\end{alerta}

\section{Bloque D --- Receptor HackRF a 1090 MHz}
\label{sec:u-receptor}

\Paso{Añada el HackRF y aplique la configuración inicial}
{\tecla{Ctrl}+\tecla{R} $\rightarrow$ dispositivo HackRF con su número de serie $\rightarrow$ \gui{OK}; ventana \textbf{HackRF Input}.}
{Frecuencia \valor{1\,090\,000 kHz}; \gui{SR} \valor{4\,000\,000 S/s}; decimación \valor{1}; \gui{Fc} \valor{Cen}; \gui{BBF} \valor{automático}; \gui{DC} \valor{activado}; \gui{IQ} \valor{activado}; \gui{Bias tee} \valor{desactivado}; \gui{RF amp} \valor{desactivado}; \gui{LNA} \valor{24 dB}; \gui{VGA} \valor{20 dB}.}
{La señal tiene un ritmo de 2 Mchip/s: el demodulador exige al menos 2 MSa/s y su documentación indica que 2 o 4 MSa/s dan buenas decodificaciones. 4 MSa/s coincide con la tasa de canal por defecto del demodulador en 7.27.2, así que no hay remuestreo. Las ganancias de partida son un punto de inicio documentado, no un óptimo: se ajustan en el bloque I con evidencia.}
{Todos los valores se aceptan.}
{Repase cada valor con su compañero contra esta lista.}
{Si la lista de dispositivos no muestra el HackRF, vuelva a la sección~\ref{sec:u-ubuntu} (permisos y \texttt{raw-usb}).}

\Paso{Arranque y examine el piso de ruido}
{Botón de arranque (triángulo azul $\rightarrow$ cuadrado verde); espectro principal.}
{Observe 2 minutos sin tocar nada.}
{Antes de buscar mensajes conviene saber cómo se ve «nada».}
{Un piso de ruido más o menos plano en 4 MHz, con el pico DC en el centro y, con tráfico, puntos o trazos muy cortos en la cascada.}
{Anote el nivel del piso (dB) y si ve trazos cortos.}
{Cuadrado rojo: reconecte el USB y arranque de nuevo.}

\begin{esperado}
En 1090 MHz no hay una portadora continua: cada mensaje dura 64 o 120~$\mu$s. El espectro instantáneo puede parecer solo ruido aunque haya tráfico; la cascada muestra, como mucho, trazos breves. Ver o no ver trazos es nivel 1: lo que decide es el demodulador.
\end{esperado}

\begin{nota}[El pico DC en 1090 MHz]
Aquí el dispositivo se centra exactamente en 1090 MHz, sobre la señal, y la corrección \gui{DC} elimina la componente continua. La señal ADS-B ocupa del orden de 2 MHz alrededor de la portadora, así que eliminar una franja muy estrecha en el centro afecta poco a la decodificación. Si lo duda, compare en el bloque I el número de tramas con \gui{DC} activado y desactivado.
\end{nota}

\section{Bloque E --- El ADS-B Demodulator}
\label{sec:u-demod}

\Paso{Añada el demodulador}
{Icono de añadir canal de la ventana \textbf{HackRF Input} $\rightarrow$ \textbf{ADS-B Demodulator} $\rightarrow$ \gui{Apply} $\rightarrow$ \gui{Close}.}
{}
{Un canal de SDRangel convierte las muestras en mensajes.}
{Ventana \textbf{ADS-B Demodulator} con controles arriba y una tabla vacía.}
{El canal aparece sobre el espectro principal, centrado en 1090 MHz.}
{Si no aparece en la lista, la instalación está incompleta: avise.}

\Paso{Revise la configuración del canal}
{Controles superiores de la ventana del demodulador.}
{Desplazamiento \valor{0 Hz}; tasa de canal \gui{SR} \valor{4M}; ancho \gui{RF BW} \valor{2,6 MHz} (por defecto); umbral de correlación \valor{7 dB} (por defecto); botón \gui{S} (\emph{Demodulate Mode S frames}) \valor{activado}; \gui{PCE} en su valor por defecto; \gui{Feed} \valor{desactivado}.}
{El umbral es el nivel en dB que debe superar la correlación con el preámbulo antes de intentar decodificar: la documentación recomienda 5--7; bajarlo produce más intentos y más tramas inválidas. \gui{S} añade tramas Mode S (ELS, EHS, MRAR) a las ADS-B. \gui{Feed} enviaría las tramas a agregadores externos: \emph{no} lo active sin autorización institucional.}
{Los valores coinciden. En 7.27.2, umbral 7 dB, Mode S activado, RF BW 2,6 MHz y tasa de canal 4 MSa/s son los valores por defecto verificados.}
{Anote en la hoja de registro todos los valores.}
{}

\Paso{Active estadísticas, gráfico y registro}
{Botones \gui{Display Statistics}, \gui{Display Chart}; registro CSV con \gui{Start/stop Logging ADS-B frames to .csv File} y \gui{.csv Log Filename}.}
{Archivo \texttt{\textasciitilde/adsb\_log\_<grupo>.csv}.}
{Las estadísticas cuentan tramas, tasa y reparto por TC y DF; el gráfico muestra tramas/s y aeronaves de los últimos 10 s; el CSV guarda cada trama (hexadecimal y correlación) como evidencia reproducible.}
{Tablas y gráfico visibles; el contador empieza en cero.}
{Tras unos minutos de tráfico, el archivo CSV crece.}
{Si no ve los botones, amplíe la ventana del demodulador.}

\begin{checkpoint}[Minuto 85: decisión de respaldo]
Espere hasta 10 minutos. ¿El total de tramas de \gui{Display Statistics} es mayor que cero y aparece al menos una fila en la tabla? \textbf{Sí}: continúe en vivo. \textbf{No}: mejore la ubicación de la antena una vez; si sigue sin tramas, pase al bloque K (sección~\ref{sec:u-offline}) y haga F--H con la grabación.
\end{checkpoint}

\section{Bloque F --- De la energía a la trama válida}
\label{sec:u-tramas}

\subsection{La escalera, aplicada a 1090 MHz}

\begin{longtable}{@{}>{\raggedright\arraybackslash}p{0.27\linewidth}>{\raggedright\arraybackslash}p{0.65\linewidth}@{}}
\toprule
\textbf{Etapa} & \textbf{Qué la demuestra} \\
\midrule
\endhead
Energía RF detectada & Trazos en la cascada; pico en el SIGLENT. Puede ser cualquier emisor de la banda. \\
Secuencia de pulsos candidata & Pulsos de 0,5~$\mu$s con la temporización de un preámbulo; el ruido también los produce por azar. \\
Preámbulo reconocido & La correlación supera el umbral (columna \emph{Correlation}). \\
Bits candidatos & 56 o 112 decisiones PPM; aún pueden tener errores. \\
Trama válida & La paridad de 24 bits es coherente con su formato (tabla~\ref{tab:df}). \\
Campos interpretados & DF, ICAO, TC y campos del mensaje decodificados correctamente. \\
Posición derivada & Cálculo CPR con un par par/impar o con una referencia (bloque H). \\
Aeronave en el mapa & Posición derivada representada; depende de la referencia configurada. \\
\bottomrule
\end{longtable}

\subsection{Preámbulo, PPM y trama}

\begin{figure}[H]
\centering
\resizebox{\linewidth}{!}{% Respuesta Mode S de enlace descendente: preámbulo de 8 us y bloque de datos PPM.
% Escala exacta: 1 us = 9 mm. Los 8 primeros bits corresponden al byte 0x8D de la
% trama real 8D4D2023... del conjunto de datos del Laboratorio 02 (DF = 17, CA = 5).
\begin{tikzpicture}[x=9mm,y=1mm,font=\sffamily\scriptsize]
  \def\h{11}
  \draw[flecha] (-0.3,0) -- (16.6,0) node[right,rotulo]{$t$ ($\mu$s)};
  \foreach \t in {0,1,...,16} \draw[MAMuted] (\t,0) -- (\t,-1.2) node[below]{\t};
  % preámbulo
  \foreach \s in {0,1,3.5,4.5} \fill[MAPrimaryDark] (\s,0) rectangle (\s+0.5,\h);
  \draw[decorate,decoration={brace,amplitude=2mm},MAPrimaryDark] (0,\h+2) -- node[above=2.5mm]{preámbulo: pulsos de 0,5 $\mu$s en 0; 1,0; 3,5 y 4,5 $\mu$s} (8,\h+2);
  % bits del byte 0x8D = 1000 1101
  \foreach \b [count=\i from 0] in {1,0,0,0,1,1,0,1} {
    \pgfmathsetmacro{\t}{8+\i}
    \ifnum\b=1 \fill[MASecondary] (\t,0) rectangle (\t+0.5,\h); \else \fill[MASecondary] (\t+0.5,0) rectangle (\t+1,\h); \fi
    \draw[MAMuted,dashed] (\t+0.5,0) -- (\t+0.5,\h+0.5);
    \node[font=\sffamily\scriptsize\bfseries] at (\t+0.5,-6.2) {\b};
  }
  \node[anchor=east] at (7.9,-6.2) {bit:};
  \draw[decorate,decoration={brace,amplitude=2mm,mirror},MASecondary] (8,-9) -- node[below=2.5mm]{DF $=10001_2=17$} (13,-9);
  \draw[decorate,decoration={brace,amplitude=2mm,mirror},MASecondary] (13,-9) -- node[below=2.5mm]{CA $=101_2=5$} (16,-9);
  \draw[decorate,decoration={brace,amplitude=2mm},MASecondary] (8,\h+2) -- node[above=2.5mm,align=center]{datos desde 8 $\mu$s: 1 bit/$\mu$s (PPM)} (16,\h+2);
  \node[rotulo,anchor=west,align=left] at (0,-22) {PPM: \textbf{1} = pulso en la \emph{primera} mitad del bit; \textbf{0} = pulso en la \emph{segunda} mitad. La línea discontinua marca la mitad.\\
     Duración total: 8 + 56 = \textbf{64 $\mu$s} (trama corta, 56 bits) u 8 + 112 = \textbf{120 $\mu$s} (trama larga, 112 bits, p.\ ej.\ DF17).};
\end{tikzpicture}
}
\caption{Respuesta Mode S de enlace descendente a escala: preámbulo de 8~$\mu$s y datos en modulación por posición de pulso (PPM). Los ocho bits mostrados son el primer byte, 0x8D, de una trama real del conjunto de datos de este laboratorio.}
\label{fig:trama}
\end{figure}

Con 4 MSa/s cada bit de 1~$\mu$s ocupa 4 muestras y cada pulso de 0,5~$\mu$s, 2 muestras: el demodulador compara la energía de las dos mitades de cada bit. Con 2 MSa/s solo hay una muestra por medio bit; por eso es el mínimo admitido.

\subsection{Formatos de enlace descendente y su validación}

\begin{longtable}{@{}>{\raggedright\arraybackslash}p{0.09\linewidth}>{\raggedright\arraybackslash}p{0.10\linewidth}>{\raggedright\arraybackslash}p{0.38\linewidth}>{\raggedright\arraybackslash}p{0.35\linewidth}@{}}
\caption{Formatos más frecuentes. La paridad usa siempre el mismo polinomio de 24 bits, pero \emph{qué} se superpone a ella cambia con el formato, y con ello lo que el receptor puede comprobar.}\label{tab:df}\\
\toprule
\textbf{DF} & \textbf{Bits} & \textbf{Contenido} & \textbf{Comprobación posible en un receptor pasivo} \\
\midrule
\endfirsthead
\toprule
\textbf{DF} & \textbf{Bits} & \textbf{Contenido} & \textbf{Comprobación posible en un receptor pasivo} \\
\midrule
\endhead
0 & 56 & Vigilancia aire-aire corta (ACAS), altitud & Paridad superpuesta con la dirección (AP): el resto \emph{es} la dirección ICAO. Solo se valida si esa dirección ya se conoce. \\
4 & 56 & Respuesta de vigilancia: altitud & AP, igual que DF0. \\
5 & 56 & Respuesta de vigilancia: código de identidad & AP, igual que DF0. \\
11 & 56 & Respuesta a \emph{all-call} y \emph{squitter} de adquisición; lleva la dirección ICAO en claro & Paridad con identificador de interrogador (II/SI): resto 0 en el squitter espontáneo, un código pequeño en respuestas a un radar. \\
16 & 112 & Vigilancia aire-aire larga (ACAS) & AP. \\
17 & 112 & \textbf{Squitter extendido ADS-B}; dirección ICAO en claro & Paridad pura: resto \textbf{0}. Es la comprobación más fuerte. \\
18 & 112 & Squitter extendido de equipos no transpondedores (p.\ ej.\ TIS-B, vehículos) & Paridad pura: resto 0. \\
20, 21 & 112 & Respuestas Comm-B (altitud o identidad + registro de 56 bits) & AP: dirección conocida necesaria. \\
\bottomrule
\end{longtable}

\begin{importante}[No todas las tramas se validan igual]
Una DF17 se valida sola: si el resto de la paridad es cero, la trama es coherente. Una DF4 o DF20 solo «se valida» si el receptor ya conoce la dirección ICAO que aparece como resto; si no, cualquier ruido produce una dirección plausible. Por eso los decodificadores aceptan estas tramas solo para aeronaves ya confirmadas por DF11 o DF17. Un pico en el espectro, por alto que sea, no es una trama válida.
\end{importante}

\subsection{Ejercicio: decodifique una trama a mano}

\begin{figure}[H]
\centering
\resizebox{\linewidth}{!}{% Campos de una trama larga DF17 (ADS-B 1090ES), 112 bits, a escala.
\begin{tikzpicture}[x=1.32mm,y=1mm,font=\sffamily\scriptsize]
  \def\campo#1#2#3#4#5{% inicio, fin, nombre, detalle, color
    \fill[#5!25] (#1,0) rectangle (#2,9); \draw[MAInk] (#1,0) rectangle (#2,9);
    \node[font=\sffamily\scriptsize\bfseries] at ({(#1+#2)/2},5.6) {#3};
    \node at ({(#1+#2)/2},2.4) {#4};}
  \campo{0}{5}{DF}{5}{MAPrimary}
  \campo{5}{8}{CA}{3}{MAPrimary}
  \campo{8}{32}{AA: dirección ICAO}{24 bits}{MASecondary}
  \campo{32}{88}{ME: mensaje ADS-B (TC = 5 primeros bits)}{56 bits}{MATertiary}
  \campo{88}{112}{PI: paridad}{24 bits}{MAWarning}
  \foreach \b in {0,5,8,32,88,112} \node[below,rotulo] at (\b,-0.5) {\b};
  \node[rotulo,anchor=west,align=left] at (0,-8) {Ejemplo verificado (conjunto de datos del Lab.\ 02): \texttt{8D 4D2023 2004D0F4CB1820 B0EFD4} $\Rightarrow$ DF 17, CA 5, ICAO 4D2023,\\TC 4 (identificación), indicativo «AMC421». La paridad de 24 bits da resto 0 (PI con II = 0).};
\end{tikzpicture}
}
\caption{Campos de una trama DF17 de 112 bits, a escala.}
\label{fig:df17}
\end{figure}

La trama \texttt{8D4D20232004D0F4CB1820B0EFD4} es una trama real de la grabación del bloque K, validada por dos decodificadores independientes y por SDRangel 7.27.2. Con la figura~\ref{fig:df17}:
\begin{enumerate}
  \item Pase \texttt{8D} a binario y obtenga DF (5 bits) y CA (3 bits).
  \item Lea la dirección ICAO (bytes 2--4).
  \item Del primer byte del campo ME (\texttt{20}) obtenga el \emph{Type Code} (5 bits) y la categoría (3 bits).
  \item Los 48 bits siguientes son 8 caracteres de 6 bits. Tradúzcalos con la tabla: 1--26 = A--Z; 32 = espacio; 48--57 = 0--9.
  \item ¿Qué tipo de mensaje es? ¿Por qué esta trama \emph{sola} no permite situar la aeronave en el mapa?
\end{enumerate}

\section{Bloque G --- La tabla de aeronaves}
\label{sec:u-tabla}

Cada fila agrupa los mensajes de una dirección ICAO. Columnas habituales en 7.27.2 (según la documentación del plugin): \emph{ICAO ID}, \emph{Callsign}, \emph{Alt}, \emph{GS}, \emph{Trk}, \emph{Hd}, \emph{VR}, \emph{Lat}, \emph{Lon}, \emph{D}, \emph{Az/El}, \emph{Frames}, \emph{ADS-B FC}, \emph{Mode S FC}, \emph{Correlation} y \emph{RSSI}. Muchas celdas aparecerán vacías: es normal.

\begin{longtable}{@{}>{\raggedright\arraybackslash}p{0.20\linewidth}>{\raggedright\arraybackslash}p{0.40\linewidth}>{\raggedright\arraybackslash}p{0.32\linewidth}@{}}
\toprule
\textbf{Campo} & \textbf{De qué mensaje sale} & \textbf{Por qué puede faltar} \\
\midrule
\endhead
ICAO & Cualquier DF11/17/18 (en claro) & Nunca falta en una fila: es la clave \\
Indicativo & DF17, TC 1--4 (identificación), menos frecuente que la posición & Aún no llegó ese mensaje; equipos sin ADS-B \\
Altitud & DF17 TC 9--18; también DF0/4/20 (Mode S) & Solo respuestas sin altitud codificada \\
Velocidad y rumbo & DF17 TC 19 & Aeronave sin ADS-B o mensaje no recibido \\
Lat/Lon & DF17 TC 9--18 (altitud barométrica) o 20--22 (altitud GNSS); en superficie, TC 5--8; siempre + cálculo CPR (bloque H) & Falta el par par/impar o una referencia válida \\
Distancia D, Az/El & Cálculo con \gui{My Position} & Derivado: erróneo si \gui{My Position} es errónea \\
Modelo, matrícula & Base de datos descargada & Nivel 7: no es RF \\
\bottomrule
\end{longtable}

\Paso{Registre tres aeronaves (o la del respaldo)}
{Tabla del demodulador.}
{Elija las filas con más tramas (\emph{Frames}).}
{Más tramas implican más campos disponibles y más confianza.}
{Filas con ICAO, y algunas con indicativo, altitud o posición.}
{Para cada valor, indique su origen: RF transmitido (nivel 5), derivado (6) o externo (7).}
{Sin filas tras 10 minutos: vea el bloque K.}

\begin{longtable}{@{}p{0.15\linewidth}p{0.12\linewidth}p{0.10\linewidth}p{0.10\linewidth}p{0.17\linewidth}p{0.08\linewidth}p{0.16\linewidth}@{}}
\caption{Registro de aeronaves (hoja de registro). Fuente: V = vivo, G = grabación.}\label{tab:u-aeronaves}\\
\toprule
\textbf{ICAO} & \textbf{Indicativo} & \textbf{Alt (ft)} & \textbf{GS (kt)} & \textbf{Lat / Lon} & \textbf{Frames} & \textbf{Fuente y niveles} \\
\midrule
\endfirsthead
\caption[]{(continuación)}\\
\toprule
\textbf{ICAO} & \textbf{Indicativo} & \textbf{Alt (ft)} & \textbf{GS (kt)} & \textbf{Lat / Lon} & \textbf{Frames} & \textbf{Fuente y niveles} \\
\midrule
\endhead
 & & & & & & \\[1.8ex]
 & & & & & & \\[1.8ex]
 & & & & & & \\[1.8ex]
\bottomrule
\end{longtable}

\begin{nota}[Ritmos de los mensajes]
La identificación, la posición y la velocidad viajan en mensajes \emph{distintos}, con ritmos distintos: la posición y la velocidad se difunden del orden de dos veces por segundo, la identificación con mucha menor frecuencia. Una aeronave puede tener altitud y no tener indicativo, o al revés, durante bastante tiempo.
\end{nota}

\section{Bloque H --- El mapa y la posición CPR}
\label{sec:u-mapa}

\subsection{Cómo se codifica la posición}

Para caber en 17 bits por coordenada, ADS-B usa \keyterm{CPR} (\emph{Compact Position Reporting}): la trama solo transmite la posición \emph{dentro de una zona} de latitud de unos 6°. Hay dos formatos, par e impar, con zonas de tamaño ligeramente distinto. La posición se reconstruye de dos maneras:
\begin{itemize}
  \item \textbf{Global:} con una trama par y una impar recientes de la misma aeronave (decenas de segundos como máximo; los decodificadores suelen exigir unos 10 s). No necesita referencia.
  \item \textbf{Local:} con una sola trama y una posición de referencia a menos de media zona (unos 180 NM). La referencia es \gui{My Position} o una posición previa de la aeronave.
\end{itemize}
Si la referencia está lejos, la decodificación local elige la zona equivocada \emph{sin avisar}. No es teoría: es lo que ocurrió al preparar este manual (figura~\ref{fig:cpr}).

\begin{figure}[H]
\centering
\resizebox{\linewidth}{!}{% Decodificación CPR local: la trama solo fija la posición DENTRO de una zona de
% latitud; la zona se elige por cercanía a la referencia (My Position).
% Valores reales medidos con SDRangel 7.27.2 sobre el conjunto de datos del Lab. 02
% (trama impar, ancho de zona 360/59 = 6,10 grados).
\begin{tikzpicture}[x=3.6mm,y=1mm,font=\sffamily\scriptsize]
  \def\dz{6.1017}
  \node[rotulo,anchor=west] at (0,25) {zonas CPR de latitud (formato impar): el punto azul es la \emph{misma} fracción codificada, repetida en cada zona};
  % referencia Bogotá y ventana válida
  \fill[MAWarningSoft] (4.7-3.05,-0.6) rectangle (4.7+3.05,10);
  \fill[MASuccessSoft] (37.0-3.05,-0.6) rectangle (37.0+3.05,10);
  \draw[flecha] (-0.5,0) -- (41,0) node[right,rotulo]{latitud (°N)};
  \foreach \k in {0,...,6} {
    \pgfmathsetmacro{\z}{\k*\dz}
    \draw[MAMuted] (\z,-1.5) -- (\z,12);
    \fill[MAPrimary] (\z+0.488,4) circle (0.55mm);
  }
  \foreach \l in {0,10,20,30,40} \node[below] at (\l,-1.6) {\l};
  \draw[MAWarning,line width=1pt] (4.7,-1) -- (4.7,10);
  \draw[MADanger,line width=1pt] (6.59,4) circle (1.1mm);
  \node[MADanger,align=center,anchor=south] at (6.59,13) {decodificado:\\$\approx$ 6,6° \textbf{(falso)}};
  \draw[MASuccess,line width=1pt] (37.0,-1) -- (37.0,10);
  \draw[MASuccess,line width=1pt] (6*\dz+0.488,4) circle (1.1mm);
  \node[MASuccess,align=center,anchor=south] at (37.1,13) {decodificado:\\37,09° (correcto)};
  \node[MAWarning,align=center,anchor=north] at (4.7,-6) {referencia automática\\$\approx$ 4,7° (Bogotá)};
  \node[MASuccess,align=center,anchor=north] at (37.0,-6) {referencia\\37,0° (zona de la grabación)};
  \node[rotulo,anchor=north,align=center] at (20.5,-16) {Banda sombreada: la referencia debe estar a menos de media zona\\($\approx$ 3° $\approx$ 180 NM) de la aeronave. Si no, el decodificador\\elige la zona equivocada \emph{sin} aviso.};
\end{tikzpicture}
}
\caption{Medido con SDRangel 7.27.2 sobre la grabación del bloque K (aeronave real cerca de Sicilia, a 37,1°N). Con la posición del receptor tomada automáticamente en Bogotá ($\approx$ 4,7°N), SDRangel situó la aeronave en Colombia, a $\approx$ 6,6°N. Con una referencia de 37,0°N, 14,0°E, la situó en 36,98--37,10°N, 13,79--13,85°E, igual que dos decodificadores independientes.}
\label{fig:cpr}
\end{figure}

\subsection{Configurar la posición del receptor}

\begin{alerta}[Privacidad de la estación]
SDRangel puede rellenar \gui{My Position} automáticamente con la ubicación del sistema operativo y, si se pulsa el botón de GPS sin datos, consultar un servicio web de geolocalización por IP. En las capturas que publique, \textbf{no muestre} las coordenadas exactas del receptor: recorte o aleje el mapa. Para los cálculos basta una precisión de 0,01° (unos 1,1 km).
\end{alerta}

\Paso{Revise y fije My Position}
{\menu{Preferences $\rightarrow$ My Position}.}
{Latitud y longitud del laboratorio en grados decimales (norte y este positivos), con dos decimales. Anote el estado de la casilla \gui{Auto-update position} (activada por defecto: toma la posición del sistema operativo).}
{Es la referencia de la decodificación CPR local y la base de la distancia \emph{D}, el acimut y el alcance máximo.}
{Los valores quedan guardados al aceptar.}
{Compare con la ubicación del aula en un mapa: diferencia menor de 0,05°.}
{Si los valores vuelven a cambiar solos, \gui{Auto-update position} está activada; para el trabajo en vivo puede dejarla si la ubicación del sistema es correcta.}

\Paso{Añada la feature Map y localice una aeronave}
{\tecla{Ctrl}+\tecla{E} (\emph{Add feature}) $\rightarrow$ \textbf{Map} $\rightarrow$ \gui{OK}. En la tabla, seleccione una fila con Lat/Lon y use \gui{Find on Feature Map}.}
{Mapa 2D. Los mosaicos del mapa se descargan de Internet: sin red, el mapa puede salir en blanco aunque las posiciones sean correctas.}
{El mapa representa objetos de varios plugins; la ADS-B aporta las aeronaves con posición calculada.}
{Símbolos de aeronave cerca de su estación; al seleccionar uno se resalta su fila.}
{La lat/lon del símbolo coincide con la de la tabla.}
{Aeronaves en la tabla pero no en el mapa: vea la tabla~\ref{tab:u-diag}.}

\begin{importante}[Mapa local frente a servicios externos]
Lo que dibuja su mapa sale de \emph{sus} tramas y \emph{su} referencia. Un servicio web de seguimiento combina miles de receptores, MLAT y otras fuentes. Si usa uno para comparar, hágalo como referencia visual, nunca como evidencia de su estación, y espere diferencias.
\end{importante}

\section{Bloque I --- Experimento de calidad de recepción}
\label{sec:u-exp}

El tráfico cambia de minuto en minuto. Para separar el efecto del receptor del efecto del tráfico, use un diseño \textbf{A--B--A}: repita la condición inicial al final. Si A y A' difieren mucho, el tráfico cambió y la comparación con B es débil.

\begin{prediccion}
¿Qué condición dará más tramas válidas por minuto? ¿Aumentará también el número de aeronaves distintas?
\end{prediccion}

{\small
\begin{longtable}{@{}p{0.07\linewidth}p{0.18\linewidth}p{0.10\linewidth}p{0.11\linewidth}p{0.11\linewidth}p{0.11\linewidth}p{0.17\linewidth}@{}}
\caption{Experimento A--B--A, 5 minutos por condición (reinicie las estadísticas con clic derecho en \gui{Display Statistics}).}\label{tab:u-exp}\\
\toprule
\textbf{Cond.} & \textbf{LNA / VGA / RF amp / umbral} & \textbf{Tramas} & \textbf{Tramas por min} & \textbf{ICAO distintos} & \textbf{Con posición} & \textbf{Correlación media / RSSI} \\
\midrule
\endfirsthead
\caption[]{(continuación)}\\
\toprule
\textbf{Cond.} & \textbf{LNA / VGA / RF amp / umbral} & \textbf{Tramas} & \textbf{Tramas por min} & \textbf{ICAO distintos} & \textbf{Con posición} & \textbf{Correlación media / RSSI} \\
\midrule
\endhead
A & 24 / 20 / off / 7 & & & & & \\[1.6ex]
B & 32 / 30 / off / 7 & & & & & \\[1.6ex]
A' & 24 / 20 / off / 7 & & & & & \\[1.6ex]
C & 24 / 20 / off / 5 & & & & & \\[1.6ex]
\bottomrule
\end{longtable}
}

\begin{itemize}
  \item Registre también: duración real, alcance máximo de \gui{Display Statistics} (solo con \gui{My Position} correcta), ubicación y orientación de la antena.
  \item \textbf{RF amp} solo como condición adicional D, con permiso del instructor, lejos de transmisores y \emph{nunca} a pocos kilómetros de un radar secundario o un DME (sección~\ref{sec:seguridad}): aporta 14 dB, pero acerca el receptor a la saturación y al límite de entrada por emisiones intensas de la banda (1030 MHz, DME, telefonía).
  \item No concluya «más ganancia es mejor» por ver más \emph{intentos}: compare tramas \emph{válidas}, número de aeronaves y estabilidad. Un umbral menor (C) aumenta las tramas pero también las inválidas.
  \item Si un día no hay tráfico, el experimento no es concluyente: declárelo así; no complete la tabla con valores supuestos.
\end{itemize}

\section{Bloque J --- El SIGLENT en 1090 MHz}
\label{sec:u-siglent}

\begin{importante}[Lo que un analizador de barrido puede y no puede ver aquí]
Una trama dura como mucho 120~$\mu$s. En modo barrido, el oscilador local pasa por 1090 MHz solo una fracción de cada barrido: la mayoría de las ráfagas llegan cuando el analizador está mirando otra frecuencia (figura~\ref{fig:barrido}a). El analizador \textbf{no decodifica} ADS-B ni permite afirmar que una ráfaga es una trama válida. Su papel es complementario: confirmar que hay actividad pulsada en 1090 MHz y estimar la duración de las ráfagas.
\end{importante}

\begin{figure}[H]
\centering
\resizebox{\linewidth}{!}{% (a) Barrido de frecuencia frente a ráfagas de 1090 MHz; (b) span cero con disparo de vídeo.
% Esquema conceptual: proporciones ilustrativas, no una medida.
\begin{tikzpicture}[x=1mm,y=1mm,font=\sffamily\scriptsize]
  \begin{scope}
    \node[anchor=west,font=\sffamily\scriptsize\bfseries] at (0,47) {(a) Barrido (span 10 MHz, RBW 1 MHz)};
    \draw[flecha] (0,0) -- (72,0);
    \node[rotulo,anchor=north east] at (72,-1) {tiempo};
    \draw[flecha] (0,0) -- (0,38);
    \node[rotulo,anchor=south west] at (1,39) {frecuencia del OL};
    \fill[MASecondarySoft] (0,16) rectangle (70,20);
    \node[rotulo,anchor=west] at (1,22.2) {1090 MHz $\pm$ RBW/2};
    \draw[MAPrimaryDark,line width=0.9pt] (2,2) -- (66,34);
    \fill[MAWarning] (30,16) rectangle (31.4,20);
    \foreach \x in {8,14,47,55,61} \fill[MASecondary] (\x,16.6) rectangle (\x+0.9,19.4);
    \draw[flechadash] (46,30) node[above,rotulo,align=center]{solo esta ventana\\observa 1090 MHz} -- (31,20.5);
    \node[rotulo,anchor=north west,align=left] at (0,-6) {Las ráfagas (barras oscuras) caen casi siempre fuera de la\\breve ventana en que el OL pasa por 1090 MHz.};
  \end{scope}
  \begin{scope}[xshift=84mm]
    \node[anchor=west,font=\sffamily\scriptsize\bfseries] at (0,47) {(b) Span cero en 1090 MHz, disparo de vídeo};
    \draw[flecha] (0,0) -- (72,0);
    \node[rotulo,anchor=north east] at (72,-1) {tiempo};
    \draw[flecha] (0,0) -- (0,38);
    \node[rotulo,anchor=south west] at (1,39) {envolvente (dB)};
    \draw[MAMuted] (0,6) -- (70,6);
    \draw[MAWarning,dashed] (0,14) -- (70,14);
    \node[MAWarning,anchor=south east] at (70,14.3) {nivel de disparo};
    \draw[MAPrimaryDark,line width=0.8pt] (0,6.5) -- (8,6.2) -- (9,25) -- (10,21) -- (11,26) -- (13,22) -- (16,26) -- (20,23) -- (26,25) -- (32,22) -- (38,26) -- (44,23) -- (49,25) -- (50,6.5) -- (70,6.2);
    \draw[<->,MAInk] (9,30) -- node[above]{$\approx$ 120 $\mu$s (trama larga)} (50,30);
    \node[rotulo,anchor=north west,align=left] at (0,-6) {Se ve la \emph{envolvente} de una trama completa; los pulsos\\de 0,5 $\mu$s y los bits no se resuelven (RBW máx.\ 1 MHz).};
  \end{scope}
\end{tikzpicture}
}
\caption{(a) Un barrido casi nunca coincide con ráfagas de microsegundos. (b) En \emph{span} cero el analizador se convierte en un receptor fijo en 1090 MHz que muestra la envolvente en función del tiempo; el disparo de vídeo congela la pantalla cuando llega una ráfaga. Esquema conceptual.}
\label{fig:barrido}
\end{figure}

Requisitos: adaptador N y una configuración de antena que en el bloque E haya dado tramas en vivo. Detenga el HackRF y traslade la antena al analizador (sección~\ref{sec:w-siglent}).

\Paso{Observación en barrido con Max Hold}
{\gui{Frequency}, \gui{Span}, \gui{BW}, \gui{Amplitude}, \gui{Trace}.}
{Centro \valor{1090 MHz}; span \valor{10 MHz}; RBW \valor{1 MHz}; VBW \valor{3 MHz}; referencia \valor{$-40$ dBm}; preamplificador \valor{activado}; traza \valor{Max Hold} durante 2 minutos.}
{Max Hold acumula las pocas coincidencias entre barrido y ráfaga.}
{Con tráfico cercano, una elevación alrededor de 1090 MHz que crece con el tiempo. Puede aparecer también actividad en 1030 MHz si cambia el centro: son interrogaciones de radares secundarios cercanos, también recibidas de forma pasiva.}
{Al volver a Clear Write, la traza instantánea parece solo ruido.}
{Sin elevación tras 2 minutos: tráfico débil o antena insuficiente; anote y pase al paso siguiente de todos modos.}

\Paso{Span cero con disparo de vídeo}
{\gui{Span $\rightarrow$ Zero Span}; \gui{Sweep $\rightarrow$ Sweep Time}; \gui{Trigger $\rightarrow$ Video}; \gui{Detect}.}
{Centro \valor{1090 MHz}; RBW \valor{1 MHz}; VBW \valor{3 MHz}; detector \valor{Pos Peak}; tiempo de barrido \valor{el mínimo que acepte el instrumento} (anote el valor); disparo de vídeo con nivel unos \valor{6--10 dB} por encima del piso de ruido.}
{En span cero el eje horizontal es tiempo; el manual del analizador indica que el disparo de vídeo solo es válido en span cero. El manual del SSA3000X Plus da un tiempo de barrido mínimo de 900~$\mu$s: con 751 puntos, cada punto representa unos 1,2~$\mu$s. Se ve la envolvente de una trama, no sus bits.}
{Al llegar una ráfaga, la pantalla se congela mostrando un bloque de unos 64 o 120~$\mu$s de ancho.}
{Use dos marcadores (\gui{Marker} en modo delta) para medir la duración del bloque.}
{Si la pantalla nunca se dispara, baje el nivel de disparo en pasos de 2 dB; si se dispara continuamente con ruido, súbalo. Si el tiempo de barrido mínimo es demasiado largo para ver la forma, declare la limitación: no es un fallo del procedimiento.}

\begin{longtable}{@{}p{0.45\linewidth}p{0.47\linewidth}@{}}
\toprule
\textbf{Observación} & \textbf{Valor} \\
\midrule
\endhead
Tiempo de barrido mínimo aceptado en span cero & \\
Nivel de disparo usado (dBm) & \\
¿Se observaron ráfagas? ¿Cuántas en 2 min? & \\
Duración medida de una ráfaga ($\mu$s) & \\
Interpretación (¿trama corta o larga?) & \\
\bottomrule
\end{longtable}

\section{Bloque K --- Respaldo offline verificado}
\label{sec:u-offline}

Este bloque es obligatorio cuando no hay recepción en vivo y recomendado siempre: permite comparar una estación real con datos conocidos. Toda evidencia obtenida aquí se etiqueta \textbf{GRABACIÓN}, nunca «en vivo».

\subsection{El conjunto de datos}

\begin{longtable}{@{}>{\raggedright\arraybackslash}p{0.27\linewidth}>{\raggedright\arraybackslash}p{0.65\linewidth}@{}}
\toprule
\textbf{Atributo} & \textbf{Valor verificado} \\
\midrule
\endhead
Archivo & {\ttfamily\nolinkurl{datos_adsb/modes1_dump1090_2MSps_1090MHz.sdriq}} (1\,427\,504 bytes) \\
SHA-256 & {\ttfamily\small\nolinkurl{d201250f486e5ff14133bc57664c8e5be6c42be3bb989529f29789674ee5f25d}} \\
Origen & \texttt{testfiles/modes1.bin} del proyecto \emph{dump1090} de Salvatore Sanfilippo (\url{https://github.com/antirez/dump1090}), añadido al repositorio el 05-01-2013 \\
Licencia & BSD de 3 cláusulas; se reproduce en \texttt{datos\_adsb/LICENSE\_dump1090.txt} \\
Formato original & I/Q de 8 bits sin signo intercalado, 2 MSa/s, centrado en 1090 MHz (formato de \texttt{rtl\_sdr}); 713\,736 bytes \\
Conversión & A \texttt{.sdriq} de SDRangel con \texttt{herramientas/adsb\_iq.py}: cabecera de 32 bytes (2\,000\,000 S/s, 1\,090\,000\,000 Hz, 16 bits, CRC-32 válido), muestras I/Q \texttt{int16} little endian $= (u - 127{,}5)\times256$ \\
Contenido & 356\,868 muestras complejas (178,4 ms); 242 tramas con paridad válida según el verificador del repositorio (DF0: 10, DF4: 4, DF5: 9, DF11: 68, DF17: 132, DF20: 13, DF21: 6), todas de la aeronave ICAO \textbf{4D2023}, indicativo \textbf{AMC421} \\
Comprobación independiente & Las 242 tramas coinciden campo a campo con la biblioteca pyModeS 3.6.0; SDRangel 7.27.2 (Windows) registró 143 tramas distintas, todas con paridad válida, 126 de ellas comunes con el verificador \\
\bottomrule
\end{longtable}

\begin{importante}[Limitaciones que debe declarar]
\begin{itemize}
  \item \textbf{Tiempo comprimido.} En 178 ms la altitud baja de 24\,275 a 20\,025 ft y la aeronave recorre unos 15 km: físicamente imposible. El archivo fue, con toda probabilidad, recortado con la opción \texttt{--snip} de dump1090, que elimina los tramos de silencio para reducir el tamaño (inferencia; la procedencia exacta no está documentada). Las tasas (tramas/s), el gráfico y los tiempos entre tramas \textbf{no tienen significado físico}.
  \item \textbf{Bucle.} SDRangel repite el archivo: los contadores crecen sin límite y la aeronave «salta» hacia atrás en cada vuelta.
  \item \textbf{Una sola aeronave}, grabada en 2013 cerca de Sicilia. La distancia \emph{D}, el acimut y el alcance máximo no significan nada respecto a su laboratorio.
  \item \textbf{Posición.} Con la referencia de su laboratorio, la posición sale \emph{falsa} (figura~\ref{fig:cpr}). Para ver la posición verdadera hay que fijar temporalmente una referencia cercana a la grabación.
\end{itemize}
\end{importante}

\subsection{Procedimiento}

\Paso{Detenga el receptor en vivo y añada FileInput}
{Detenga \textbf{HackRF Input}. \tecla{Ctrl}+\tecla{R} $\rightarrow$ \textbf{FileInput} $\rightarrow$ \gui{OK}.}
{}
{\textbf{FileInput} sustituye al HackRF como fuente de muestras: el resto de la cadena no cambia.}
{Una ventana \textbf{FileInput} con botón de apertura de archivo, reproducción/pausa, bucle y aceleración.}
{}
{}

\Paso{Abra la grabación}
{Botón de abrir archivo de \textbf{FileInput} (deshabilitado mientras reproduce).}
{\texttt{datos\_adsb/modes1\_dump1090\_2MSps\_1090MHz.sdriq}; bucle \valor{activado}; aceleración \valor{1}. En Ubuntu, el SDRangel instalado como snap lee con seguridad archivos de su carpeta personal: copie allí la carpeta \texttt{datos\_adsb}.}
{La cabecera aporta tasa y frecuencia: no se teclean.}
{La ventana muestra \valor{2000 kS/s}, \valor{1\,090\,000 kHz}, duración \valor{00:00:00.178} y fecha de inicio del 05-01-2013 (hora local).}
{Esos cuatro valores coinciden. SDRangel rechazaría el archivo si el CRC de la cabecera no fuera válido.}
{Si no acepta el archivo, compruebe su SHA-256 (sección~\ref{sec:u-sha}).}

\Paso{Fije temporalmente una referencia cercana a la grabación}
{\menu{Preferences $\rightarrow$ My Position}.}
{Anote sus valores actuales. Desactive la casilla \gui{Auto-update position}. Introduzca \valor{latitud 37,0}, \valor{longitud 14,0}.}
{La decodificación local necesita una referencia a menos de unos 180 NM de la aeronave.}
{Los valores se mantienen al reabrir el diálogo.}
{Si vuelven solos a otra posición, \gui{Auto-update position} sigue activada: desactívela.}
{Sin poder cambiarla: realice el paso y documente que la posición del mapa es \emph{falsa}; esa observación es en sí un resultado válido sobre CPR.}

\Paso{Añada el demodulador y reproduzca}
{Icono de añadir canal de \textbf{FileInput} $\rightarrow$ \textbf{ADS-B Demodulator}. Botón de reproducción.}
{Desplazamiento \valor{0 Hz}; tasa de canal \gui{SR} \valor{2M} (igual a la del archivo); umbral \valor{7 dB}; \gui{S} activado.}
{Con 2M no hay remuestreo de un archivo de 2 MSa/s. Esta configuración es la que se ejecutó al verificar el respaldo.}
{En segundos aparece una fila \textbf{4d2023}, indicativo \textbf{AMC421}.}
{Compare con la tabla~\ref{tab:u-referencia}.}
{Sin filas tras 30 s: compruebe que la reproducción avanza y que el canal está en 0 Hz.}

\begin{longtable}{@{}>{\raggedright\arraybackslash}p{0.34\linewidth}>{\raggedright\arraybackslash}p{0.58\linewidth}@{}}
\caption{Resultados de referencia (grabación, no en vivo), observados al verificar el respaldo con SDRangel 7.27.2 y con el verificador del repositorio.}\label{tab:u-referencia}\\
\toprule
\textbf{Campo} & \textbf{Rango esperado} \\
\midrule
\endfirsthead
\caption[]{(continuación)}\\
\toprule
\textbf{Campo} & \textbf{Rango esperado} \\
\midrule
\endhead
ICAO / indicativo & 4D2023 / AMC421 \\
Altitud & 20\,025--24\,275 ft (varía en cada vuelta del bucle) \\
Velocidad respecto al suelo / rumbo & 371--390 kt / $\approx$ 158° \\
Posición (referencia 37,0 N, 14,0 E) & 36,96--37,10 °N; 13,79--13,86 °E \\
Formatos presentes & DF0, 4, 5, 11, 17, 20, 21 \\
Correlación de preámbulo registrada & $\approx$ 6--63 dB (mediana $\approx$ 23 dB) \\
\bottomrule
\end{longtable}

\Paso{Restaure su posición}
{\menu{Preferences $\rightarrow$ My Position}.}
{Los valores anotados y el estado anterior de \gui{Auto-update position}.}
{Olvidarlo falsearía distancias y posiciones en vivo en las sesiones siguientes.}
{}
{Reabra el diálogo y compruebe los valores.}
{}

\subsection{Comprobación independiente del archivo}
\label{sec:u-sha}

En cualquier sistema con Python 3 y \texttt{numpy}, desde la carpeta \texttt{lab\_SDR\_surveillance\_LaTeX}:

\Comando{sha256sum datos_adsb/modes1_dump1090_2MSps_1090MHz.sdriq}
{Comprobar que el archivo no se alteró.}
{\texttt{d201250f486e5ff1...74ee5f25d} (valor completo en la tabla del conjunto de datos).}
{Otro valor: archivo dañado o modificado.}
{Vuelva a descargarlo del repositorio.}

\Comando{python3 herramientas/adsb_iq.py verificar datos_adsb/modes1_dump1090_2MSps_1090MHz.sdriq}
{Demodular el archivo con un decodificador independiente y simple, sin corrección de errores.}
{\texttt{178.434 ms, 430 preámbulos candidatos, 242 tramas válidas, 1 direcciones ICAO} y la lista de tramas con su regla de validación.}
{\texttt{ModuleNotFoundError: numpy}.}
{\texttt{python3 -m pip install numpy} en un entorno virtual, o use el entorno del repositorio.}

\begin{esperado}
Observe en la salida del verificador que 430 \emph{preámbulos candidatos} producen 242 tramas válidas: los demás son coincidencias con ruido o tramas con errores. Ese cociente ilustra por qué un preámbulo no es todavía una trama. El verificador también marca 4 posiciones «sospechosas», calculadas con pares par/impar separados en el tiempo real: una lección directa sobre CPR.
\end{esperado}

\section{Bloque L --- Informe de la estación}
\label{sec:u-informe}

\begin{entregable}[Laboratorio 02]
\begin{enumerate}
  \item \textbf{Hardware:} conector, ubicación y cable de la antena; decisión de preparación (\PASS, \COND\ o \BLOQ) y su justificación.
  \item \textbf{Software:} \texttt{lsb\_release -ds}, versión de SDRangel (\menu{Help $\rightarrow$ About}) y salida de \texttt{hackrf\_info}.
  \item \textbf{Receptor:} todos los parámetros de \textbf{HackRF Input} y del \textbf{ADS-B Demodulator}.
  \item \textbf{Señal:} captura del espectro y la cascada en 1090 MHz (nivel 1) y, si existe, del SIGLENT.
  \item \textbf{Tramas:} captura de \gui{Display Statistics} y las primeras 20 líneas del CSV; la trama decodificada a mano (bloque F).
  \item \textbf{Tabla de aeronaves} (tabla~\ref{tab:u-aeronaves}) con el origen de cada valor (niveles 5, 6 o 7).
  \item \textbf{Mapa:} captura \emph{sin} coordenadas exactas del receptor.
  \item \textbf{Estadísticas:} tabla~\ref{tab:u-exp} completa o declarada no concluyente.
  \item \textbf{Origen de la evidencia:} cada captura etiquetada EN VIVO o GRABACIÓN.
  \item \textbf{Limitaciones y conclusiones} (máximo una página): qué demostró su estación, qué no, y qué cambiaría primero para mejorarla.
\end{enumerate}
\end{entregable}

\section{Solución de problemas}
\label{sec:u-diag2}

\begin{longtable}{@{}>{\raggedright\arraybackslash}p{0.24\linewidth}>{\raggedright\arraybackslash}p{0.33\linewidth}>{\raggedright\arraybackslash}p{0.35\linewidth}@{}}
\caption{Diagnóstico del Laboratorio 02.}\label{tab:u-diag}\\
\toprule
\textbf{Síntoma} & \textbf{Diagnóstico} & \textbf{Acción} \\
\midrule
\endfirsthead
\toprule
\textbf{Síntoma} & \textbf{Diagnóstico} & \textbf{Acción} \\
\midrule
\endhead
\texttt{lsusb} no muestra 1d50:6089 & Cable, puerto o alimentación & Otro cable de datos y puerto directo; \texttt{sudo dmesg | tail} al reconectar. \\
\texttt{hackrf\_info} solo funciona con \texttt{sudo} & Regla udev o grupo & Bloque C: leer la regla, \texttt{groups}, \texttt{usermod}, cerrar sesión. \\
SDRangel no lista el HackRF & Snap sin \texttt{raw-usb} u otro programa usando el HackRF & \texttt{snap connections sdrangel}; cerrar \texttt{hackrf\_info} y otros programas; reconectar USB. \\
Botón de arranque en rojo & Error del dispositivo & Detener, reconectar USB, arrancar. \\
Espectro con ruido y sin trazos & Tráfico escaso, antena o ubicación & Antena vertical en una ventana; cable C; esperar 10 min; si nada, bloque K. \\
Ruido muy alto o espurios & Ganancia excesiva o interferencia del PC & LNA 24 / VGA 20, RF amp off; alejar la antena del PC y de cargadores. \\
Muchos intentos, ninguna fila & Umbral muy bajo o señales débiles & Umbral a 7 dB; mejore la ubicación; no baje de 5 dB. \\
Filas sin indicativo & Mensajes de identificación aún no recibidos & Esperar; es normal. \\
Tabla con filas pero mapa vacío & Sin posición CPR, sin mosaicos o \gui{My Position} errónea & Esperar pares de posición; comprobar Internet para mosaicos; revisar \gui{My Position}. \\
Aeronave en un lugar imposible & Referencia CPR lejana (figura~\ref{fig:cpr}) & Corrija \gui{My Position}; en la grabación, use 37,0 N / 14,0 E. \\
Distancias o alcance absurdos & \gui{My Position} errónea o datos de la grabación & Corrija la posición; en la grabación, D y alcance no tienen sentido. \\
FileInput no abre el archivo & Archivo dañado o reproducción en curso & Pausa; comprobar SHA-256; volver a descargar. \\
Interfaz lenta & CPU saturada por mapa 3D o muchas trayectorias & Mapa 2D; cerrar ventanas que no use. \\
SIGLENT no muestra nada en 1090 & Barrido frente a pulsos (figura~\ref{fig:barrido}) & Max Hold 2 min o span cero con disparo de vídeo; declarar la limitación. \\
\bottomrule
\end{longtable}

\section{Cuestionario final}
\label{sec:u-final}

Individual. Respuestas y justificaciones en la guía del instructor.

\CuestionarioFinalU
